\section{正则量子化：相空间、预量子化与极化}
\label{sec:canonical}

第 \ref{sec:cs} 节只写了作用量。这一节回答：为什么“沿曲面切开”能得到一个有限维的量子 Hilbert 空间？它由哪几步构成？每一步的数学前提是什么？

\subsection{第一步：切开流形，读出辛形式}

把三维流形写成 $M=\Sigma\times\mathbb{R}$，取 $A_0=0$ 的规范。作用量 \eqref{eq:cs} 中对时间是一阶的，三阶项因只剩两个空间指标而消失 \src{367}：
\begin{equation}
L=\frac{k}{8\pi}\int \dd t\int_\Sigma \epsilon^{ij}\Tr\Bigl(A_i\frac{\partial}{\partial t}A_j\Bigr).
\label{eq:sympl}
\end{equation}
这已经是“$p\dot q$”的形式，所以 $A_i$ 与 $A_j$ 是共轭变量，泊松括号直接读出 \src{367}：
\begin{equation}
\{A_i^a(x),A_j^b(y)\}=\frac{4\pi}{k}\,\epsilon_{ij}\,\delta^{ab}\,\delta^2(x-y).
\label{eq:poisson}
\end{equation}

\subsection{第二步：高斯定律与辛商}

$A_0$ 是拉格朗日乘子，变分给出高斯定律约束 \src{367}：
\begin{equation}
\epsilon^{ij}F^a_{ij}=0\quad\Longleftrightarrow\quad F=0 .
\label{eq:gauss}
\end{equation}
二维曲面只有一个独立分量，所以约束就是平坦性。它是第一类约束，生成规范变换。

经典上分两步施加：先限制到 $F=0$（平坦联络），再按规范群取商。这一步就是 Marsden--Weinstein 辛商，用的是 Atiyah--Bott 的框架 \src{367}（文献 [33]）。结果是
\begin{equation}
\cM=\Hom(\pi_1\Sigma,G)/G ,
\end{equation}
即“$\Sigma$ 上平坦联络模去规范变换”。

\begin{keypoint}
$\cM$ 是\textbf{有限维}的：$g>1$ 时 $\dim_\mathbb{R}\cM=(2g-2)d$，$d=\dim G$。它紧、带奇点、辛体积有限。半经典地看，相空间每单位体积容纳一个量子态，所以量子 Hilbert 空间有限维 \src{368}。这正是本文“可精确求解”的技术前提——如果相空间是无限维的，后面所有有限维论证都不成立。
\end{keypoint}

\subsection{第三步：预量子化}

几何量子化的标准套路是：先预量子化，再选极化。

\textbf{预量子化的前提是整性}：$[\omega]\in H^2(\cM,\ZZ)$，否则不存在预量子化线丛。Witten 在 §3.1 里\textbf{没有讨论这一条}——他直接取 $\cM_J$ 上的行列式线丛 $\cL$（带 Quillen 度规，文献 [34]），并写 $[\omega]=c_1(\cL^{\otimes k})$ \src{368}。这不是疏漏，而是构造性绕过 \ours：行列式线丛由几何本身给出，它的第一陈类自动是整的，所以“整性条件”在这个具体问题里是免费的。真正需要整性论证的是一般辛流形，不是这里。

于是预量子化 Hilbert 空间是 $\cL^{\otimes k}$ 的截面空间；但截面空间太大（同时依赖位置与动量），必须再选极化。

\subsection{第四步：极化与全纯截面}

紧相空间上不能用 $q,p$ 分离的实极化（$p_j$ 无界），所以 Witten 取 \textbf{K\"ahler 极化}：选 $\cM$ 上的复结构 $J$，把 $\cM_J$ 看成 K\"ahler 簇，量子态空间取全纯截面 \src{368}：
\begin{equation}
\cH_\Sigma^{(J)}=H^0\bigl(\cM_J,\cL^{\otimes k}\bigr).
\label{eq:hilbert}
\end{equation}
这是有限维的。

\begin{keypoint}
\textbf{极化在数学上做了什么}：它把“函数”换成“沿极化方向协变常数的截面”。实极化下这是沿拉格朗日叶的常截面；K\"ahler 极化下这是全纯截面。李思第 5.4 节把这一步讲得比本文细：先证明 $\cL$ 上的联络与辛形式相容（$[\nabla_X,\nabla_Y]=\nabla_{[X,Y]}-\ii\hbar\,\omega(X,Y)$），再定义极化，再取极化截面，最后才得到量子态空间。本文跳过了相容性这一层，直接使用结论。
\end{keypoint}

\subsection{第五步：$J$ 变时发生了什么}

$\cH_\Sigma^{(J)}$ 全纯依赖于 $J$。Witten 的断言是：这些空间拼成复结构模空间上的一个\textbf{射影平坦}向量丛 \src{369--371}；“与 $J$ 无关”的准确含义是存在典范的平坦联络，而由于共形反常，平坦只能到射影意义 \src{370}。这个联络就是 Hitchin 联络（带标记点时是 KZ 联络）；Segal 后来严格证明了它的平坦性 \src{370}（文献 [16]）。

单值性有物理意义：绕模空间一周，态空间上的变换是映射类群的作用 \src{371}。这正是二维共形场论里“模函子”的雏形。

\subsection{第六步：被略去的半形式修正}

几何量子化一般还需要半形式修正（取 $K^{1/2}$，$K$ 是典范线丛）。本文 §3 略去了它，只用 $\cL^{\otimes k}$ \src{369}。

后果是能级平移：量子化给出的组合是 $k+h^\vee$（$h^\vee$ 为对偶 Coxeter 数；$SU(N)$ 时 $h^\vee=N$），最终配分函数里的量子参数是 \src{382}
\begin{equation}
q=\exp\frac{2\pi\ii}{k+N}\qquad(SU(N)),\qquad
q=\exp\frac{2\pi\ii}{k+2}\qquad(SU(2)).
\label{eq:q}
\end{equation}
读本文时若直接用 $q=e^{2\pi\ii/k}$，会与 §4 的最终公式差一个平移。

\begin{keypoint}
\textbf{两处必须分清}：本文的平移其实是在 §2 以单圈/$\eta$ 修正的形式出现的——(2.16)(2.17) 把耦合写成 $k+c_2(G)/2$，Witten 自己指出这与二维流代数里“$k\to k+c_2/2$”是同一件事 \src{359}。半形式修正只是几何量子化语言里对同一现象的另一种描述。若按标准 $K^{1/2}$ 去推，$K_{\cM}=\mathrm{Det}^{-2h^\vee}$，给出的可能是 $k-h^\vee$——符号取决于 $\omega$ 与 $\mathrm{Det}$ 的取向约定，本文没有给。所以“$k\to k+h^\vee$”应理解为“有效组合是 $k+h^\vee$”，不要当作一条可独立推导的定理。
\end{keypoint}

\begin{boundary}
\begin{itemize}[leftmargin=1.4em]
\item 原文给出：(3.1)--(3.3)、辛商、相空间有限维、预量子化线丛、K\"ahler 极化、射影平坦性与单值性。
\item 我们补的：把“整性免费”解释为行列式线丛的构造性绕过（S2-A 的判断）；把极化步骤与李思 5.4 节对齐。
\item 未证：半形式修正导致 $k\to k+h^\vee$ 的独立推导（本文未给，三本教材也都没有讲 $K^{1/2}$ 机制，只给等价陈述）；射影平坦性的严格证明归 Segal。
\end{itemize}
\end{boundary}
